\subsubsection{Warum nur ein Schatten auch nach beliebig langer Zeit
nicht
reicht}\label{warum-nur-ein-schatten-auch-nach-beliebig-langer-zeit-nicht-reicht}

Für alle t gilt exakt \(\operatorname{tr}[P(t)h]=0\). Der Operatororbit
einer einzigen Paarsonde hat auf dem Zweizustandsfaktor Rang drei statt
vier. Mit sämtlichen inneren Sonden bleibt der unsichtbare Raum

\[\adjustbox{max width=\linewidth}{$\displaystyle 
h\otimes\operatorname{Herm}(60).
$}\]

Selbst bekannte Gesamtspur eins lässt 3599 unsichtbare Richtungen übrig.
Man benötigt eine Information über die zweite Zweigmatrix, nicht bloß
noch mehr Zeitpunkte desselben unvollständigen Bildes. Bei a=0 oder Δ=0
verliert die angegebene Drei-Zeit-Rekonstruktion weitere Richtungen. Das
sind genaue Grenzen, keine mangelnde numerische Genauigkeit.

\subsection{7. Echter kausaler Eingriff ohne zusätzlichen
Hoppingterm}\label{echter-kausaler-eingriff-ohne-zusuxe4tzlichen-hoppingterm}

Jede aktive W-Spalte gehört genau einer Zeile; jede Zeile enthält acht
disjunkte Fermionpaare. Der ganze N=2-Raum zerfällt deshalb in 60
invariante neundimensionale Sterne und 1536 ungekoppelte Paare.
Einschließlich der sieben dunklen Richtungen je Stern ergibt sich wieder
dunkle Dimension 1956.

Im ersten Stern starte man in \(f_4^\dagger f_{57}^\dagger|0\rangle\).
Nach nativer Vorentwicklung wird entweder nichts getan oder die gezielte
lokale Phase \(Z_4=(-1)^{n_4}\) angewendet. Nach weiterer identischer
Entwicklung wird \(n_5\) ausgelesen:

\[\adjustbox{max width=\linewidth}{$\displaystyle 
|4,57\rangle\ \xrightarrow{U_T}\
\{I\text{ oder }Z_4\}\ \xrightarrow{U_T}\ n_5.
$}\]

Z4 erhält N, ist eine unitäre spurtreue Operation und kommutiert mit n5.
Unmittelbar beim Eingriff bleibt die Besetzungsstatistik von Mode 5
deshalb für jeden Zustand gleich. Der spätere Effekt ist somit kein
direkter Zugriff auf den Empfänger.

Setze

\[\adjustbox{max width=\linewidth}{$\displaystyle 
\Omega=\sqrt{\Delta^2/4+8g^2},\quad T=\pi/\Omega,
\quad d=\Delta/(2\Omega),\quad x=4\cos^2(\pi d/2).
$}\]

In beiden Armen gilt am Ende exakt Nf=2 und Nb=0, ohne Nachselektion.
Die Wahrscheinlichkeiten und ihre Differenz sind

\[\adjustbox{max width=\linewidth}{$\displaystyle 
p_{\rm frei}=\frac{x(4-x)}{64},\quad
p_{\rm Impuls}=\frac{9x^2}{1024},\quad
\boxed{\delta_{4\to5}=\frac{x(25x-64)}{1024}.}
$}\]

Für \(0<g^2\le3\Delta^2/32\) ist die Differenz strikt negativ. Bei Δ=1
und g=1/20 erhält man numerisch

\[\adjustbox{max width=\linewidth}{$\displaystyle 
T\approx6.0459978807807,\quad
p_{\rm frei}\approx0.00087491599417248,\quad
p_{\rm Impuls}\approx0.0000017344871305812.
$}\]

Das entspricht etwa 0.08749 Prozent gegenüber 0.0001734 Prozent
Besetzungswahrscheinlichkeit. Der Effekt ist klein, aber der
Nichtnullnachweis ist exakt. Ein unabhängiger numerischer Lauf auf dem
gesamten 2076-dimensionalen Hamiltonoperator bestätigt die vollständigen
Zustände und die Negativkontrolle ohne Vorentwicklung.

\textbf{Geschlossen ist die bedingte mathematische Frage:} Der
ursprüngliche H kann eine nachweisbare Moden-zu-Moden-Wirkung
vermitteln; ein zusätzlicher Linkterm ist für diesen Zeugen nicht nötig.
\textbf{Nicht geschlossen ist die Herkunft} der Produktpräparation, der
einzeln adressierbaren Phase und des Besetzungsinstruments. Der
Phasenimpuls ist eine hinzugewährte Kontrolle, auch wenn die freie
Dynamik unverändert bleibt. Die Moden sind noch keine nachgewiesenen
räumlich getrennten Labore. Dies ist außerdem nicht die Ausbreitung des
früher bewiesenen N=64-Entnahmepols.

\subsubsection{7.1 Noch kleiner: derselbe Dreizustandsraum trägt
Schatten und
Eingriff}\label{noch-kleiner-derselbe-dreizustandsraum-truxe4gt-schatten-und-eingriff}

Für den gerade beschriebenen Versuch genügt sogar der gemeinsame Raum

\[\adjustbox{max width=\linewidth}{$\displaystyle 
\mathcal K_3=\operatorname{span}\{|p_0\rangle,|R_7\rangle,|b_0\rangle\},
\quad |R_7\rangle=\tfrac1{\sqrt7}\sum_{q=1}^7s_q|p_q\rangle.
$}\]

Die s sind die ursprünglichen W-Vorzeichen. Dieser Raum enthält die
notwendige dunkle Richtung; er wird nicht mit einem einzelnen hellen
Zweizustandsblock verwechselt. Mit der expliziten Isometrie J gilt ohne
Rest HJ=Jh3 und Z4J=Jz3, wobei

\[\adjustbox{max width=\linewidth}{$\displaystyle 
h_3=\begin{pmatrix}0&0&-g\\0&0&\sqrt7g\\-g&\sqrt7g&\Delta\end{pmatrix},
\quad z_3=\operatorname{diag}(-1,1,1),\quad
e_3=J^\dagger n_5J=\operatorname{diag}(0,1/7,0),
\quad b_3=\operatorname{diag}(0,0,1).
$}\]

Definiere L(O)=i{[}h3,O{]} und C(O)=z3Oz3. Die neun hermiteschen
Operatoren

\[\adjustbox{max width=\linewidth}{$\displaystyle 
e_3,b_3,Le_3,Lb_3,L^2e_3,L^2b_3,L^3e_3,CL^2e_3,CL^2b_3
$}\]

besitzen in den angegebenen reellen Matrixkoordinaten die Determinante

\[\adjustbox{max width=\linewidth}{$\displaystyle 
\boxed{\det M=8\Delta^3g^{10}/343\ne0.}
$}\]

Sie bestimmen folglich jeden Zustand auf genau dem Träger des kausalen
Versuchs. Die benötigten Daten sind terminale
Besetzungswahrscheinlichkeiten und ihre zeitlichen Ableitungen bis
Ordnung drei, mit und ohne vorangestellten Phasenimpuls. Bei bekannter
Normierung genügt sogar die autonome Zeitfamilie mit Ableitungen bis
Ordnung vier. Der Impuls ist für diesen normierten Tomographiesatz also
nicht zwingend, wohl aber der untersuchte Eingriff im kausalen
Vergleich.

Dies verbindet die beiden Aufgaben auf derselben kleinen Ausführung,
nicht nur auf Teilräumen mit ähnlichen Spektren. Es setzt bekannte
Dynamik, identische frische Zustände, kalibrierte Zeiten und präzise
Ableitungsdaten voraus. Ein endliches robustes Abtastverfahren ist noch
nicht konstruiert. Die im Teilbericht exakt gerechnete konservative
Fehlergrenze kann bei schwacher Kopplung große Fehlerverstärkung
zulassen.

\textbf{Wichtige Instrumentengrenze:} n5 erhält den Raum nicht. Es gilt

\[\adjustbox{max width=\linewidth}{$\displaystyle 
(n_5J-Je_3)^\dagger(n_5J-Je_3)=\operatorname{diag}(0,6/49,0).
$}\]

Eine nichtselektive projektive n5-Messung lässt aus R7 das Gewicht 12/49
außerhalb von K3 zurück. Die Reduktion gilt bis zur Endmessung und für
Tomographie mit frischen Exemplaren, nicht für beliebige weitere
Messfolgen desselben Exemplars. Ein wirklich grobes Lüders-Instrument
auf alle sieben Restpaare würde den Raum erhalten; feines Auslesen mit
anschließendem Vergessen der Labels ist ein anderes Instrument und
verliert sogar Gewicht 6/7. Gerade diese Unterscheidung muss ein
ursprüngliches Aufzeichnungsinstrument respektieren.

\subsection{8. Was Holografie hier sinnvoll
bedeutet}\label{was-holografie-hier-sinnvoll-bedeutet}

Bildlich sind die verschiedenen Beschreibungen Ansichten derselben
Maschine. Eine Ansicht zeigt Besetzungen, eine andere innere
Paarstruktur, eine dritte das spätere Verhalten nach einem Eingriff.
Ihre gemeinsame Kalibrierung kann die verborgenen Phasen bestimmen.
Abschnitt 6 realisiert genau diese Idee im endlichen hellen Raum.

Das ist zunächst Informationsrekonstruktion, nicht bereits gravitative
Holografie. Es wurden weder eine physische Raumgrenze noch ein
Flächengesetz oder eine Bulk-Rand-Dualität hergeleitet. Das Wort
Schatten bekommt hier einen überprüfbaren Inhalt: eine konkrete
Operatorabbildung, ihren Kern, ein gezieltes zusätzliches Messergebnis
und eine zurückgehaltene Vorhersage. Ähnliche Zahlen oder Spektren ohne
gemeinsame Abbildung genügen weiterhin nicht.

\subsection{9. Nächste entscheidende
Nachweise}\label{nuxe4chste-entscheidende-nachweise}

\textbf{Übergeordnete Frage: Welche ursprüngliche Kompositionsregel
erzeugt gemeinsam unterscheidbare Vorkommen, ihre Bewegung, ihren
Zustand und ihre Instrumente?} Die letzte nachgereichte fundamentale
Reduktion trifft diese Frage. Ein genaueres lokales Spektrum kann sie
nicht ersetzen. Die folgenden kleinen Prüfungen dienen als überprüfbare
Annahmekriterien für eine solche Regel; sie sind keine Ersatzdefinition
von Raum und kein Beweis, dass es zwingend getrennte kopierte Banken
sein müssen.

\textbf{Erste Priorität: den kleinen Versuch an den Compiler
anschließen.} Ein einziger vollständiger Vertrag für Präparation, Phase
und terminale Auslesung wäre aussagekräftiger als eine weitere
allgemeine Matrixfamilie. Als Ziel dienen der konkrete N=2-Eingriff und
sein abgesicherter Effekt. Zu liefern ist ein ursprüngliches Instrument
oder ein expliziter Adapter aus dem vorhandenen Kontextträger
einschließlich Zustand, Ladung und Aufzeichnung. Ein Matrixeintrag oder
eine Gruppenzugehörigkeit allein besteht diesen Test nicht. Eine
begründete globale Zustandsregel darf dabei nicht mit einer externen
Laborpräparation des ganzen Universums verwechselt werden.

\textbf{Zweite Priorität: dieselbe niedrige geladene Anregung bewegen.}
Die neuen kleinen Sektoren lösen den früheren N=64-Anschluss nicht
automatisch. Gesucht ist ein nichtskalarer Übergang auf der tatsächlich
abgesicherten geladenen Antwort, mit kontrollierter Restkopplung und auf
demselben Grundzustand. Reine Überlappung innerhalb eines exakt
entarteten Polraums erzeugt weiter kein Hopping. Matrixmomente und die
geprüfte Blockresolvente sind dafür ein konkreter Ansatz.

\textbf{Dritte Priorität: den Feldadapter einschließlich Boosts
schließen.} Der neue Händigkeitstest unterscheidet einen gültigen
gleichhändigen Projektor von einem nur unter Drehungen gültigen Adapter.
Die konkrete geladene Mode muss mit richtigen Indizes, CAR,
W-Kontraktion und ihrem Matrixelement auf Ω abgebildet werden. Parallel
kann ein aus verfügbaren Operationen oder dynamischen Freiheitsgraden
hervorgehendes P(λ) mit kontrollierter Lücke und nichtverschwindender
Krümmung untersucht werden. Die reale Zweikomponentenfamilie allein
liefert keinen solchen nativen Geometrieschritt. Feld- und
Geometrieprüfung müssen an derselben Konstruktion bleiben.

Die separaten mathematischen Fragen dieser Runde werden damit
beantwortet, ohne einen Abschluss der übergeordneten Physik zu
behaupten.

\subsection{10. T1 bis T8 nach dieser
Runde}\label{t1-bis-t8-nach-dieser-runde}

Die Bezeichnungen folgen der Arbeitszuordnung der geprüften Synthese;
die jeweiligen vollständigen Tore bleiben offen.

{\def\LTcaptype{none} % do not increment counter
\begin{longtable}[]{@{}
  >{\raggedright\arraybackslash}p{(\linewidth - 4\tabcolsep) * \real{0.3333}}
  >{\raggedright\arraybackslash}p{(\linewidth - 4\tabcolsep) * \real{0.3333}}
  >{\raggedright\arraybackslash}p{(\linewidth - 4\tabcolsep) * \real{0.3333}}@{}}
\toprule\noalign{}
\begin{minipage}[b]{\linewidth}\raggedright
Tor
\end{minipage} & \begin{minipage}[b]{\linewidth}\raggedright
Neuer Beitrag
\end{minipage} & \begin{minipage}[b]{\linewidth}\raggedright
Noch fehlender entscheidender Inhalt
\end{minipage} \\
\midrule\noalign{}
\endhead
\bottomrule\noalign{}
\endlastfoot
T1 Ursprung/Auswahl & Tatsächlicher W-/Clock-Replay, vollständiger
N-Erhaltungsausschluss des dokumentierten Alphabets & Auswahl des
Prozesskerns und ausführbarer Primitive; physisches g/Δ \\
T2 Half-Charge/E8 & Neue exakte Operator- und Ressourcenbilanz &
Renormierter geladener Adapter mit Energie-, Adjungierten-, E8- und
Clockkontrolle \\
T3 gemeinsamer 3+1D-Ursprung & Kausaler Eingriff unter ursprünglichem H
im festen N=2-Sektor & Abgeleitete lokale Teile, Distanz, Raumdimension
und gemeinsamer 3+1D-Parent \\
T4 chirale Materie & Exakter gleichhändiger Ausleseprojektor und
gegengängiger Boost-Gegenzeuge & Nativer relativistischer Feldadapter,
chirales Maß, Anomalien, Spiegelentkopplung \\
T5 Kontinuum/Dynamik & Exakte endliche Schließungen, bedingte
Mehrzeitrekonstruktion & Größenuniforme Fehler, relativistischer
wechselwirkender Limes, Streuung \\
T6 Parameter/Spektren & Neue ν5, strengere Ritz-/Polschranke und
korrigierte geladene Variationsmomente & Physische Kopplungsauswahl und
einheitliche Übertragung zu gemessenen Größen \\
T7 Gravitation & Projektorgeometrie bestätigt, native reale
Berry-Zweiform eingegrenzt & Positiver dynamischer Spin 2 und
universelle Kopplung aus demselben Parent \\
T8 Zustand/Instrumente & Referenzverbrauch, explizite Messressourcen und
Phasenkontrolle bilanziert & Ursprüngliche Präparation, Apparate,
Aufzeichnung und Wiederverwendung \\
\end{longtable}
}

Es gibt auch keinen neuen Beweis von RH, effizienter allgemeiner
Faktorisierung oder einer P-versus-NP-Aussage. Diese Runde testet eine
konkrete TFPT-Prozesskette und überträgt endliche Ergebnisse nicht auf
andere offene Probleme.

\subsection{11. Reproduktionsumfang}\label{reproduktionsumfang}

Der gemeinsame Replay führt alle eigenen mathematischen Prüfer normal
und optimiert mit Warnungen als Fehler aus. Die Teilberichte
unterscheiden exakte Identitäten von numerischen Kontrollen. Die 57
Prüfbedingungen des gelieferten Papers und dessen 20 übernommene
Kettenprüfungen werden separat ausgewiesen, nicht als neue eigene
Entdeckungen addiert. Sie decken nicht automatisch jeden Satz des Papers
ab; die Korrektur in §8.4 zeigt dies ausdrücklich.

Das Paket enthält sämtliche unveränderten Eingaben, Teilberichte und
Ergebnisdateien. Der Extraktions-Replay verwendet nur eingefrorene
Dateien, keine sich verändernden Repositoryquellen. Die große
ursprüngliche Grundzustandsprüfung und die gesamte frühere 822er-Suite
wurden in dieser Runde nicht erneut ausgeführt. Frühere entsprechende
Befunde werden als Vorbefunde bewahrt, nicht als frisch bestätigt
ausgegeben.

\begin{center}\rule{0.5\linewidth}{0.5pt}\end{center}

\section{Nachgereichte Texte: neue Norm, korrigierte Zustandszuordnung
und
Feldadapter}\label{nachgereichte-texte-neue-norm-korrigierte-zustandszuordnung-und-feldadapter}

Die drei während der laufenden Prüfung nachgereichten Texte wurden
vollständig gelesen und unverändert eingefroren. Sie sind keine
Anweisungen, Ergebnisse ungeprüft zu übernehmen. Zwei Texte vereinfachen
ältere Arbeitsstände stark; die dritte längere Forschungsnotiz enthält
ausführlichere Beweise einer bereits bekannten grundsätzlichen Diagnose.

\subsection{A. Ein tatsächlich neuer Rechenwert trägt eine stärkere
Schranke}\label{a.-ein-tatsuxe4chlich-neuer-rechenwert-truxe4gt-eine-stuxe4rkere-schranke}

Die Norm der fünften reinen Bosonstufe ist gegenüber dem hiesigen
v1.6.7-Snapshot neu:

\[\adjustbox{max width=\linewidth}{$\displaystyle 
\nu_5=1\,866\,738\,327\,552\,000.
$}\]

Die eingefrorene ursprüngliche Spur-/Wick-Rechnung wurde unabhängig
erneut ausgeführt: alle Normen von ν1 bis ν5, bei ν5 genau 840
Wick-Netzwerke in 93 vorzeichentreuen Isomorphieklassen, mit
Python-Ganzzahlen statt gerundeter Netzwerkarithmetik. Normaler und
optimierter Lauf liefern identische Ergebnisbytes. Dies ist ein frischer
Replay derselben Methode, keine zweite unabhängige Methode für ν5. Die
kleineren Normen stimmen mit den bereits vorhandenen unabhängigen
Zertifikaten überein.

Aus den sechs normierten Vektoren v0 bis v5 entsteht eine echte
Variationskompression. Ihre Diagonale ist kΔ, ihre benachbarte Kopplung
\(g\sqrt{\nu_{k+1}/\nu_k}\). Fügt man den bereits exakt konstruierten
orthogonalen Zwei-Boson-Rücklaufrest w2 hinzu, ergibt sich eine
siebendimensionale Kompression. Seine zusätzliche Kopplung liegt zu v3;
sie ist \(g\|w_2\|/\sqrt{\nu_3}\). Die übrigen angegebenen Nullstellen
folgen aus Bosonzahl und Orthogonalität.

Unabhängige rationale Sturm-Zählung liefert

\[\adjustbox{max width=\linewidth}{$\displaystyle 
-1.13847420<E_{6,\rm Ritz}/\Delta<-1.13847419,
$}\] \[\adjustbox{max width=\linewidth}{$\displaystyle 
-1.13847610<E_{7,\rm Ritz}/\Delta<-1.13847609.
$}\]

Dies sind \textbf{keine invarianten Unterräume} und nicht die sechste
bzw. siebte Stufe der vollständigen Hamilton-Lanczoskette. Sie sind
dennoch gültige Variationsräume. Deshalb folgt streng

\[\adjustbox{max width=\linewidth}{$\displaystyle 
\boxed{E_0<-1.13847609\Delta.}
$}\]

Zusammen mit dem beibehaltenen, in dieser Runde nicht vollständig neu
bewiesenen N=63-Satz \(E_h>-1.121899\Delta\) erhält man die verbesserte
Entnahmegrenze

\[\adjustbox{max width=\linewidth}{$\displaystyle 
\boxed{\epsilon=E_h-E_0>0.01657709\Delta.}
$}\]

Sie ersetzt im aktuellen Stand die schwächere untere Schranke
0.00773911Δ aus v1.6.7. Die obere Grenze 0.039079764Δ und das bisherige
Liniengewicht über 88.007628 Prozent werden dadurch nicht automatisch
verändert. Es wurde kein exakter Zentralwert des Grundzustands oder Pols
gefunden.

